% ============================================================
% Avant-propos
% Présentation des outils utilisés pour la production du rapport :
% LaTeX, GitHub et GitHub Actions.
% ============================================================

\sectionspeciale{Avant-propos}

\subsection*{Sur le choix de \LaTeX{}}

Le présent rapport est rédigé en \LaTeX{}, un système de composition de
documents largement utilisé dans le milieu universitaire et l'édition
scientifique. J'ai eu l'occasion de me familiariser avec cet outil lors
de la rédaction de mon mémoire de Master~2 en génie civil, dont les
sources sont consultables à l'adresse suivante :
\url{https://drive.google.com/drive/folders/1xuWxaYyATofsG46Qz3gQ8edoHxM41sDl?usp=drive_link}.

Le choix de \LaTeX{} pour ce rapport tient à plusieurs raisons concrètes.
D'abord, la qualité typographique des formules mathématiques et le soin
apporté à la mise en page des équations longues, qui sont nombreuses
dans un rapport de calcul de structures. Ensuite, la stabilité des
références croisées : chaque figure, chaque tableau, chaque équation
porte une étiquette interne, et leur numérotation se met à jour
automatiquement au fil des modifications, sans risque de
désynchronisation entre le texte et les renvois. Enfin, le format de
sortie est un document PDF reproductible, c'est-à-dire qu'à partir des
mêmes fichiers source, n'importe quel lecteur obtient le même PDF
identique au caractère près.

Le rapport est organisé en plusieurs fichiers source : un fichier
principal qui orchestre la compilation, un préambule contenant les
réglages de mise en page, et un fichier par séance regroupant le texte,
les calculs et les figures de la séance. Cette organisation modulaire
facilite la relecture et permet de travailler sur une séance sans
toucher aux autres.

\subsection*{Sur l'usage de GitHub}

Les fichiers source du rapport sont hébergés sur GitHub, une plateforme
publique de stockage de code et de documents textuels. L'adresse du
dépôt est :
\url{https://github.com/fatoumata-gc/ossature-batiment-chec}.

GitHub joue ici le rôle d'archive consultable et historique. Chaque
modification apportée à un fichier est enregistrée avec sa date, son
auteur et une description, ce qui permet de retrouver à tout moment
l'état du rapport à une étape antérieure du travail. Cette traçabilité
n'est pas une fin en soi : elle reflète simplement la même rigueur
qu'un cahier de laboratoire daté et paraphé. Pour un examinateur qui
souhaiterait consulter directement les sources, vérifier une formule
dans son contexte de calcul ou observer la progression du travail au
cours du semestre, l'historique du dépôt en offre la possibilité.

\subsection*{Sur GitHub Actions, et ce que ce n'est pas}

Une seule fonctionnalité de GitHub mérite ici une explication un peu
plus précise, car son nom peut prêter à confusion : il s'agit de
GitHub~Actions.

GitHub~Actions est un service d'automatisation. Concrètement, il permet
de demander au serveur d'exécuter une tâche programmée chaque fois
qu'un fichier source est modifié. Dans le cas du présent rapport, la
tâche est unique et très simple : à chaque mise à jour du dépôt, le
serveur lance la commande \texttt{pdflatex} sur le fichier
\texttt{main.tex}, recompile le document et publie le PDF résultant.
C'est exactement ce que je ferais en local avec un Makefile ou un
script shell, à la différence près que la tâche est exécutée par le
serveur de GitHub plutôt que sur ma propre machine. L'intérêt pratique
est qu'on dispose toujours d'une version PDF à jour sans devoir penser
à recompiler manuellement après chaque correction.

Il me paraît important de préciser ce que GitHub~Actions \emph{n'est
pas}, car le terme « actions » peut évoquer des fonctionnalités
beaucoup plus ambitieuses qu'il ne recouvre.

GitHub~Actions n'est pas un système d'intelligence artificielle. Il ne
génère aucun contenu, ne propose aucune formulation, ne corrige aucune
phrase, ne retouche aucune équation. Le service ne fait que rejouer une
suite d'instructions définies à l'avance dans un fichier de
configuration : récupérer les fichiers source, installer la
distribution \LaTeX{}, exécuter \texttt{pdflatex} en deux passes,
produire un PDF, et le rendre téléchargeable. C'est un déclencheur de
compilation, équivalent moderne d'un Makefile. La rédaction du rapport,
le choix des hypothèses, le déroulé des calculs et la formulation des
conclusions relèvent intégralement du travail de l'auteur.

\subsection*{Lecture du rapport}

Le rapport peut être lu indépendamment de toute considération
informatique : le PDF est un document autonome qui se suffit à
lui-même. Les éléments présentés dans cet avant-propos visent
uniquement à donner au lecteur la possibilité, s'il le souhaite, de
remonter aux sources du document, de consulter son historique, ou de
recompiler le PDF à partir du dépôt public.

\clearpage
